\subsection{The cost of the service mesh (RQ4, H3)}
\label{sec:res-mesh}

Table~\ref{tab:mesh} lists the mesh's cost at every rate where the comparison is clean. Per hop, the two
sidecars added \xo{mesh/browse/200/perhop}--\xo{mesh/checkout/25/perhop}\,ms of transit, of the same order as
the 0.18--0.25\,ms (P90--P99) that Istio reports for a sidecar pair in version 1.22~\cite{Istio122Perf}, and the database
round trips of a request took \xo{mesh/browse/200/db}--\xo{mesh/checkout/200/db}\,ms longer. The mesh explained
\xo{mesh/browse/50/share}\% of D's mean end-to-end latency on browse at the lowest rate and
\xo{mesh/checkout/25/share}\% on checkout. The share rose with load, to \xo{mesh/browse/400/share}\% and
\xo{mesh/checkout/200/share}\%. The difference D$-$C itself stayed nearly constant; the share rose because D's own
latency fell with load, as the host's low-load penalty receded.

H3 requires the mesh to account for at least a quarter of end-to-end latency on the six-hop path and its per-hop
cost not to shrink as the path gets deeper. On checkout at the lowest rate, where the condition is evaluated, the
mesh accounted for \xo{H3/share_deep}\% of end-to-end latency (median over all clean rates:
\xo{H3/share_deep_median}\%). At each workload's lowest rate the per-hop transit cost was
\xo{H3/per_hop_shallow_ms}\,ms on the three-hop browse path and \xo{H3/per_hop_deep_ms}\,ms on the six-hop
checkout path. It did not shrink as the path got longer, so the total grows at least linearly with the number of
hops. The two
per-hop intervals overlap, so the data show linear growth, not compounding. \textbf{H3 is supported.} Our original
statement of H3 also predicted that the mesh would become the single largest contributor above four hops. On
checkout at the lowest rate the mesh's \xo{mesh/checkout/25/e2e}\,ms is indeed larger than any single layer of
configuration~C (whose largest, database time, is \xo{layer/checkout/C/25/db}\,ms).

The sidecars' CPU is the larger cost. At \xo{ref/browse/high} browse requests per second they consumed
\xo{cpumsref/browse/D/mesh}\,ms of CPU per request, more than the whole application tier
(\xo{cpumsref/browse/D/app}\,ms); for checkout at \xo{ref/checkout/high} requests per second,
\xo{cpumsref/checkout/D/mesh} against \xo{cpumsref/checkout/D/app}\,ms. Because the sidecars' CPU lies outside the
equal application budget, D's capacity is not lower than C's; under a budget that included them, it would very
likely be lower.
